\section{Invariant Distributions: General Balance and Detailed Balance}
\label{sec:A.2}

In this section, we discuss the notion of statistical equilibrium, which is central to the theory of
Markov chains.

\begin{definition}[General Balance Equation; Invariant Distribution]
\label{def:A.7}
A Markov kernel $p(x, z)$ satisfies the general balance equation with respect to $\pi$ if
\[
  \pi(x) = \int_E \pi(z)\, p(z, x)\, dz.
\]
We then say that $\pi$ is an invariant distribution of the Markov kernel $p(x, z)$.
\end{definition}

Note that if $E$ is discrete, the general balance equation is simply $\pi = \pi P$.

\begin{definition}[Ergodic Markov Chain]
\label{def:A.8}
A Markov chain $\mathbb{X}$ is called ergodic if there exists a distribution $\pi$ such that, for
all $A \subset E$ and initial distributions $\pi_0$, it holds that
\[
  \lim_{n \to \infty} \Prob(X_n \in A) = \pi(A).
\]
We say that $\pi$ is the limit distribution of $\mathbb{X}$.
\end{definition}

\begin{lemma}
\label{lem:A.9}
Let $\mathbb{X}$ be an ergodic Markov chain with Markov kernel $p(x, z)$ and limit distribution
$\pi$.  Then, $\pi$ is an invariant distribution for $p(x, z)$.  In particular, if $\mathbb{X}$ is
initialized at statistical equilibrium ($X_0 \sim \pi$), then $X_n \sim \pi$ for all $n \geq 0$.
\end{lemma}

\begin{proof}
We only give a sketch of the proof.  By the dominated convergence theorem
\begin{align*}
  \lim_{n \to \infty} \pi_{n+1}(x)
    &= \lim_{n \to \infty} \int_E \pi_n(z)\, p(z, x)\, dz \\
    &= \int_E \lim_{n \to \infty} \pi_n(z)\, p(z, x)\, dz,
\end{align*}
which shows that $\pi(x) = \int_E \pi(z)\, p(z, x)\, dz$.  The final claim of the lemma is proved
by induction.
\end{proof}

The general balance equation implies that the transition probabilities of $\mathbb{X}$ preserve
equilibrium.  For this reason, we call a distribution $\pi$ that satisfies the general balance
equation~(\ref{def:A.7}) an invariant distribution of the chain $\mathbb{X}$.  If $\mathbb{X}$ is
ergodic, we can use the general balance equation to show that $\pi$ is the equilibrium distribution
of the chain.  However, given a distribution $\pi$, it is often difficult to find a Markov kernel
for which $\pi$ satisfies the general balance equation.  It is actually easier to find a Markov
kernel that satisfies a \emph{stronger} condition, known as detailed balance.

\begin{definition}[Detailed Balance]
\label{def:A.10}
We say that a transition density $p(x, z)$ satisfies detailed balance with respect to $\pi$ if,
for all $x, z \in E$, it holds that $\pi(x)\, p(x, z) = \pi(z)\, p(z, x)$.
\end{definition}

The Metropolis Hastings algorithm in Chapter~4 provides a general recipe to define a family of
kernels that satisfy detailed balance with respect to a given distribution.  The following result
shows that detailed balance implies general balance.

\begin{theorem}[Detailed Balance Implies General Balance]
\label{thm:A.11}
Let $p(x, z)$ be a Markov kernel that satisfies detailed balance with respect to a distribution
$\pi$.  Then, $\pi$ is an invariant distribution for $p(x, z)$.
\end{theorem}

\begin{proof}
Using the assumption of detailed balance and that $p(x, z)$ is a Markov kernel (and so it
integrates to one over its second argument) gives
\[
  \int_E \pi(z)\, p(z, x)\, dz
  = \int_E \pi(x)\, p(x, z)\, dz
  = \pi(x) \int_E p(x, z)\, dz
  = \pi(x). \qedhere
\]
\end{proof}

\begin{remark}
\label{rem:A.12}
Detailed balance is not \emph{necessary} for general balance.  Detailed balance implies that the
chain is time reversible: in equilibrium the chain behaves the same whether it is run forward or
backward in time.  However, there are many ergodic Markov chains that are not time reversible.
\end{remark}
